\section{Qualitative Analysis}
\label{sec:qualitative}

The aggregate metrics show that silent tool poisoning is a recurring failure mode. We inspect representative trajectories to clarify what the metrics capture. The examples below are selected from rescored result files and illustrate mechanisms behind the aggregate trends.

\subsection{Blind Compliance Looks Like Successful Analysis}

The most direct failure mode is copying a poisoned observation into the final answer without further checks. In a multi-table customer-spend task, the clean joined rows identify the correct segment spend, but the poisoned observation swaps the reported entity binding. The base LangGraph ReAct adapter adopts the poisoned result and restates a plausible calculation. Nothing in the final answer indicates that the agent recognized the observation as untrustworthy. This is the core BCR failure: the tool call succeeds, the output looks analytical, and the agent treats that successful observation as reliable evidence.

The same pattern appears in semantic tasks. In a treatment/control label-flip case, the poisoned observation states that \texttt{Control} has the highest conversion rate while preserving the clean rate \texttt{0.145}. The base adapter returns \texttt{Control -- 0.145}. The number is plausible, but the entity binding is wrong. This illustrates why ToxicBench evaluates more than arithmetic error: a data agent must preserve the link between values, labels, rows, and evidence sources.

\subsection{Validation Must Recompute Evidence}

Successful recovery requires evidence, not just cautionary language. In the same join task family, the guarded LangGraph ReAct adapter performs an independent verification pass, re-inspects the relevant tables, recomputes the join, and returns the clean entity--value binding. In the treatment/control task, it re-inspects the dataframe and returns \texttt{Treatment -- 0.145}. These trajectories explain why the full guard drives BCR to zero in both the numerical and semantic/schema ablations: the final answer is grounded in a second observation after the first corrupted observation.

This behavior differs from simply adding a caution instruction. The useful signal is not that the agent says tool outputs may be unreliable; it is that the agent obtains fresh evidence and bases the final answer on the recomputation. This motivates reporting VR and RR separately from ADR. Detection without correction is not enough for data-analysis workflows.

\subsection{Budgeted Guards Can Under-Recover}

The light guard is informative as a budget ablation. It also reduces observed BCR to zero and remains close to the full guard in aggregate, but its tighter verification budget can still under-recover on individual tasks. In the expanded numerical suite, some failures remain ordinary analysis or answer-extraction errors rather than blind copying: the guarded run identifies that verification is needed, but the final answer still misses the expected numeric tolerance. These cases show that cheaper validation can prevent direct copying while still failing to recover the correct answer.

\subsection{Implications}

The cases support three conclusions. First, high clean-task success is insufficient because poisoned observations can look like normal analytical evidence. Second, semantic poisoning exposes evidence-binding failures that are not captured by numeric tolerance alone. Third, guarded verification works when it forces an independent check, but validation budget matters; too little budget changes blind compliance into under-recovery rather than solving the task.
